%% ==================== 问题 X 建模思路与数据准备 模板 (v1.5.6, BZD 数模社 2026 规范) ====================
%%
%% 借鉴 BZD 数模社 第五章 模板: 5.X.1 建模思路 + 数据处理
%%
%% BZD 结构 (5.X.1):
%%   - 建模思路: 一段话说明任务 + 数据特征 + 数学类型 + 模型选择理由
%%   - 数据处理: 数据来源 + 处理方法 + 处理后保留样本 + 特征构造
%%   - 算法选择: 模型/算法名称 + 选择依据
%%
%% 跑题用户必须:
%% 1. 删掉所有 % 注释
%% 2. 替换【】占位符为自己的实际内容
%% 3. 删 下面的 \ModelingSlot{...} 行 (check 15 抓这个 inline 【】)

\subsubsection{问题 X 建模思路与数据准备}

\textbf{（1）建模思路} \quad 问题 X 要求【问题目标, 1 句。例如: 分析影响 Y 染色体浓度稳定性的关键因素】。根据题目条件和数据特征, 该问题在数学上属于【数学类型。例如: 影响因素识别 + 多元回归分析】。考虑到【数据特点, 1 句。例如: 特征维度高 (10 个变量) 且变量间可能存在共线性】, 选择【主模型, 如 多元线性回归 (MLR) + 随机森林回归 (RFR) 组合】进行求解, 并针对【本题特点】对模型的【目标函数/约束/参数/算法】进行调整。

\textbf{（2）数据处理} \quad 本问题使用【数据来源, 1 句。例如: 题目附件 Sheet1 男胎 1082 例】。首先进行【处理方法, 1 句。例如: 缺失值检测 + 异常值剔除 (3$\sigma$ 准则) + 标准化 (z-score)】, 原因是【处理理由, 1 句。例如: 保证数据质量并消除量纲影响】。处理后保留【样本数量】个有效样本, 并构造【特征名称, 1 句】作为模型输入。

\textbf{（3）算法选择与依据} \quad 为保证模型可解释性, 同时具备非线性捕捉能力, 本文采用【主模型 + 备选模型】组合方案:
\begin{itemize}
  \item \textbf{主模型}:【如: 多元线性回归 (MLR), 数学类型为线性回归, 用于识别主要线性影响因素】
  \item \textbf{备选模型}:【如: 随机森林回归 (RFR), 集成 100 棵决策树, 用于捕捉非线性关系并提供特征重要性排序】
  \item \textbf{选择依据}:【如: MLR 可解释性强便于系数分析, RFR 拟合优度通常更高 (R²) 且提供特征重要性, 两者对比可避免单一模型的局限】
\end{itemize}

\textbf{（4）模型评价指标} \quad 本文采用以下指标评价模型表现:
\begin{itemize}
  \item $R^2$ (决定系数): 衡量模型解释因变量变异的能力, 越接近 1 越好
  \item $\text{RMSE}$ (均方根误差): 衡量预测值与真实值的偏差, 越小越好
  \item $\text{MAE}$ (平均绝对误差): 衡量预测值与真实值的绝对偏差, 越小越好
\end{itemize}

% 跑题后必须删下面这一行 (check 15 占位符自检)
\textbf{【这里粘贴建模思路 + 数据处理 + 算法选择 拼起来的不超过 1 页内容, 删掉这一行】}
